\KOMAoption{paper}{a4, portrait}
\areaset{267mm}{180mm}
\fancyheadoffset{0pt} % обновляем хедер и футер

\newgeometry{
  top=14mm,
  bottom=11mm,
  right=10mm,
  left=22mm,
}

\begin{landscape}

%\setlength{\headsep}{7.5mm} 
%\setlength{\topmargin}{-27mm}
\setlength{\footskip}{15.5pt}

\color{unimain}{\section[(рекомендуемое) Подключение дискретных входов \newline и выходов устройства к внутренним цепям алгоритмов для архитектуры I типа]{(рекомендуемое)\\Подключение дискретных входов и выходов устройства к внутренним цепям алгоритмов для архитектуры I типа}\label{app:inouts}}
\color{black}

\begin{figure}[h!]
  \centering
  
\includegraphics[width=1.1\textwidth]{\apppath/Приложение. Подкл вх и вых/_latex/img/1.pdf}
  \caption{Подключение дискретных входов и выходов к внутренним сигналам алгоритмов для блока в слоте М1}
  \label{fig:sig1} % Опционально: ссылка на рисунок
\end{figure}

\begin{figure}[h!]
  \centering
  \includegraphics[width=1.5\textwidth]{\apppath/Приложение. Подкл вх и вых/_latex/img/2.pdf}
  \caption{Подключение дискретных входов и выходов к внутренним сигналам алгоритмов для блоков в слотах М3-М5}
  \label{fig:sig2} % Опционально: ссылка на рисунок
\end{figure}

\begin{figure}[h!]
  \centering
  \includegraphics[width=1.5\textwidth]{\apppath/Приложение. Подкл вх и вых/_latex/img/3.pdf}
  \caption{Подключение дискретных входов и выходов к внутренним сигналам алгоритмов для блоков в слотах М6-М8}
  \label{fig:sig3} % Опционально: ссылка на рисунок
\end{figure}

\begin{figure}[h!]
  \centering
  \includegraphics[width=1.5\textwidth]{\apppath/Приложение. Подкл вх и вых/_latex/img/4.pdf}
  \caption{Подключение дискретных входов и выходов к внутренним сигналам алгоритмов для блоков в слотах М9-М11}
  \label{fig:sig4} % Опционально: ссылка на рисунок
\end{figure}

\clearpage
\end{landscape}
\restoregeometry


